\chapter{Fonds propres réglementaires pour les portefeuilles de négociation}\label{ch:rc:regulatory-capital-for-trading-books}

Au titre des règles de risque de marché que le Comité de Bâle a finalisées en janvier 2019, une banque qui
souhaite fixer les fonds propres d'un desk au moyen de son propre modèle doit démontrer, chaque trimestre,
que le modèle explique le P\&L propre du desk. Si le P\&L du modèle de risque et celui des systèmes de valorisation
divergent trop, le desk paie une majoration~; s'ils divergent davantage, le desk perd le modèle
et se replie sur une formule standardisée qui, pour le portefeuille de ce chapitre, est presque trois fois plus élevée.
Les règles ont mis des années à entrer en vigueur~: reportées une première fois par le Comité, repoussées de deux ans
dans l'Union européenne, retardées au Royaume-Uni, et encore à l'état de projet aux États-Unis en 2026. Ce chapitre
calcule les deux approches sur le portefeuille du chapitre~21~: les sensibilités standardisées,
l'\omterm{def:rc:market-risk-measures:es}{expected shortfall} avec sa calibration en période de tensions,
et le test d'attribution qui tranche entre les deux.

\section{Le portefeuille de négociation et sa frontière}

\begin{definition}[Fundamental Review of the Trading Book]\label{def:rc:regulatory-capital-for-trading-books:frtb}
Le \emph{Fundamental Review of the Trading Book}\index{Fundamental Review of the Trading Book}
(FRTB) est le cadre de fonds propres au titre du risque de marché établi par le Comité de Bâle en janvier 2019~: une frontière révisée
entre les portefeuilles, une approche standardisée fondée sur les sensibilités, et une \omterm{def:rc:regulatory-capital-for-trading-books:ima}{approche des modèles internes}
fondée sur l'\omterm{def:rc:market-risk-measures:es}{expected shortfall}, assortie de tests d'éligibilité au niveau du desk.
\end{definition}

\begin{definition}[Portefeuille de négociation, portefeuille bancaire]\label{def:rc:regulatory-capital-for-trading-books:books}
Le \emph{portefeuille de négociation}\index{portefeuille de négociation} regroupe les instruments détenus pour une revente à court terme, pour
tirer profit de mouvements de prix à court terme, pour verrouiller des profits d'arbitrage, ou pour couvrir
de telles positions, valorisés à la juste valeur chaque jour par le P\&L~; tout le reste relève du \emph{portefeuille
bancaire}\index{portefeuille bancaire}, dont le risque de marché (principalement le risque de taux) est géré comme au
chapitre~24.
\end{definition}

La frontière compte, car les deux portefeuilles supportent des fonds propres différents~; transférer des instruments
de l'un à l'autre pour réduire les fonds propres est restreint et, lorsque cela est permis, tout avantage en fonds propres est
écarté.

\section{L'approche standardisée~: sensibilités, courbure, défaut et risque résiduel}

\begin{definition}[Méthode fondée sur les sensibilités]\label{def:rc:regulatory-capital-for-trading-books:sbm}
La \emph{méthode fondée sur les sensibilités}\index{méthode fondée sur les sensibilités} calcule des sensibilités prescrites
(un PV01 par maturité pour les taux d'intérêt, la variation de valeur pour un mouvement de 1~\% pour le change),
les multiplie par des pondérations de risque prescrites, et agrège les sensibilités pondérées au sein
des compartiments ($K_b = \sqrt{\max(0,\sum_{kl}\rho_{kl}WS_kWS_l)}$) puis entre les compartiments
($\sqrt{\sum_bK_b^2+\sum_{b\ne c}\gamma_{bc}S_bS_c}$), selon trois scénarios de corrélation (les corrélations prescrites,
1.25 fois ces corrélations plafonnées à un, et $\max(2\rho-1,0.75\rho)$)~; l'exigence est le
plus grand des trois totaux.
\end{definition}

\begin{definition}[Exigence au titre du risque de courbure]\label{def:rc:regulatory-capital-for-trading-books:curv}
L'\emph{exigence au titre du risque de courbure}\index{exigence au titre du risque de courbure} capture les pertes des options
au-delà de leur delta~: chaque \omterm{def:rc:market-risk-measures:factor}{facteur de risque} est choqué à la hausse et à la baisse de sa pondération de risque, et la perte
au-delà du terme de delta, $\mathrm{CVR}^\pm = -\sum_i(V_i(x^\pm)-V_i(x)\mp\mathrm{RW}\,s_i)$, est retenue, la plus
grande des deux directions.
\end{definition}

\begin{definition}[Exigence au titre du risque de défaut, majoration pour risque résiduel]\label{def:rc:regulatory-capital-for-trading-books:drc}
L'\emph{exigence au titre du risque de défaut}\index{exigence au titre du risque de défaut} capture le \omterm{def:rc:reduced-form-credit:jtd}{risque de saut au défaut} des
positions de crédit et d'actions que les chocs de spread ne voient pas. La \emph{majoration pour risque
résiduel}\index{majoration pour risque résiduel} applique une part forfaitaire du notionnel aux risques que les
sensibilités ne capturent pas~: 1~\% pour les instruments aux sous-jacents exotiques, 0.1~\% pour les autres
risques résiduels.
\end{definition}

\omcode{code/firm/frtb/firm_frtb.py}{37}{63}{Agrégation au sein des compartiments et entre les compartiments, et exigence GIRR au titre du delta sous un scénario de corrélation.}\label{lst:rc:regulatory-capital-for-trading-books:sbm}

\begin{example}[L'exigence standardisée du portefeuille]\label{ex:rc:regulatory-capital-for-trading-books:sa}
Portefeuille du chapitre~21, au 23 septembre 2026. Ses sensibilités GIRR (qui ramènent les sensibilités au taux au pair
des deux obligations aux maturités de deux et dix ans en dollars, une simplification) valent $+56\,525$
et $-155\,012$ dollars par point de base~; avec des pondérations de risque de $1.3~\%/\sqrt2$ et $1.1~\%/\sqrt2$ et une
corrélation de 88.7~\%, l'exigence au titre du delta GIRR s'élève à 7.83~millions de dollars. Le delta de change (EUR et JPY, pondération de risque
$15~\%/\sqrt2$, corrélation entre devises de 60~\%) s'élève à 5.30~millions de dollars~; le vega de change, soit le vega du straddle multiplié par
sa volatilité de 8~\% à une pondération de risque de 100~\%, à 14.49~millions de dollars~; la courbure de change, obtenue en choquant EURUSD de
$\pm10.6~\%$, à 39.23~millions de dollars. Dans le scénario de faible corrélation, le total s'élève à 68.29~millions de dollars, le
plus grand des trois (\cref{fig:rc:regulatory-capital-for-trading-books:sa}).
\end{example}

\begin{omfigure}
\begin{tikzpicture}
\begin{axis}[width=11.5cm,height=5.4cm,ybar,bar width=8pt,
  symbolic x coords={GIRR delta,FX delta,FX vega,FX curvature},xtick=data,xticklabels={{GIRR delta},{FX delta},{FX vega},{FX courbure}},
  ylabel={charge (millions de dollars)},
  tick label style={font=\small},label style={font=\small},
  ymin=0,ymax=45,grid=major,grid style={gray!25},
  legend columns=3,legend style={font=\footnotesize,at={(0.5,-0.2)},anchor=north,draw=none,fill=none,/tikz/every even column/.append style={column sep=8pt}},
  table/col sep=comma]
\addplot[fill=omMeth!60,draw=omMeth] table[x=part,y=low]{figdata/rates-credit-risk/23-regulatory-capital-for-trading-books/sa.csv};
\addplot[fill=omDef!60,draw=omDef] table[x=part,y=medium]{figdata/rates-credit-risk/23-regulatory-capital-for-trading-books/sa.csv};
\addplot[fill=omThm!60,draw=omThm] table[x=part,y=high]{figdata/rates-credit-risk/23-regulatory-capital-for-trading-books/sa.csv};
\legend{corrélation faible,moyenne,forte}
\end{axis}
\end{tikzpicture}

\omcaption{Charges standardisées du portefeuille du chapitre~21, par composante et par scénario de corrélation.
La courbure du straddle court domine~; les charges de delta varient selon le scénario (des corrélations plus faibles
suppriment la compensation entre les positions à deux et dix ans), tandis que le vega et un \omterm{def:rc:rates-risk:factors}{facteur de courbure} unique ne varient pas.
Données~: le tutoriel du chapitre.}\label{fig:rc:regulatory-capital-for-trading-books:sa}
\end{omfigure}

\section{L'approche des modèles internes~: expected shortfall et horizons de liquidité}

\begin{definition}[Approche des modèles internes, horizon de liquidité]\label{def:rc:regulatory-capital-for-trading-books:ima}
L'\emph{approche des modèles internes}\index{approche des modèles internes} fixe les fonds propres à partir du
modèle propre de la banque d'\omterm{def:rc:market-risk-measures:es}{expected shortfall} à 97.5~\%, calculé sur un horizon de base de dix jours et majoré
pour les \omterm{def:rc:market-risk-measures:factor}{facteurs de risque} dont les \emph{horizons de liquidité}\index{horizon de liquidité} sont plus longs (10, 20, 40,
60 ou 120 jours, durée supposée nécessaire pour sortir d'une position ou la couvrir sans faire bouger les prix)~:
$\mathrm{ES} = \sqrt{\mathrm{ES}_T(P)^2+\sum_{j\ge2}\bigl(\mathrm{ES}_T(P,j)\sqrt{(\mathrm{LH}_j-\mathrm{LH}_{j-1})/T}\bigr)^2}$.
\end{definition}

\begin{definition}[Expected shortfall en période de tensions]\label{def:rc:regulatory-capital-for-trading-books:ses}
L'\emph{expected shortfall en période de tensions}\index{expected shortfall en période de tensions} calibre l'ES sur la
période de douze mois la plus sévère pour le portefeuille de la banque, calculée sur un ensemble réduit de \omterm{def:rc:market-risk-measures:factor}{facteurs de risque}
dotés d'un long historique, et majorée par le rapport entre l'ES actuel avec tous les facteurs et
l'ES actuel avec l'ensemble réduit (au moins un).
\end{definition}

L'idée est plus ancienne que la FRTB~: les révisions du Comité de Bâle de juillet 2009 (Bâle~2.5) ont ajouté à l'exigence au titre de la VaR
une VaR en période de tensions, soit une VaR à dix jours au seuil de 99~\% du portefeuille actuel, calibrée sur une période continue
de douze mois de tensions importantes, comme 2007--08, chaque exigence étant retenue au plus haut entre sa dernière valeur et un
multiple de sa moyenne sur soixante jours.

Le capital pour les facteurs modélisables vaut $\mathrm{IMCC} = \rho\,\mathrm{ES}(\text{tous})+(1-\rho)\sum_i\mathrm{ES}(\text{classe }i)$ avec $\rho = 0.5$,
soit la moitié de l'ES diversifié et la moitié de la somme sur les classes de risque, et l'exigence est la plus élevée entre
l'IMCC de la veille et 1.5 fois sa moyenne sur soixante jours, auquel s'ajoute la charge en période de tensions pour les
facteurs non modélisables.

\begin{example}[L'exigence du modèle interne]\label{ex:rc:regulatory-capital-for-trading-books:ima}
Sur des variations à dix jours qui se chevauchent, l'ES à 97.5~\% du portefeuille sur la dernière année est de 5.27~millions de dollars. Parcourant
toutes les fenêtres de douze mois depuis 2016, la plus sévère se termine le 3 octobre 2022
(la \cref{fig:rc:regulatory-capital-for-trading-books:windows})~: son ES vaut 14.41~millions de dollars pour l'ensemble du portefeuille,
5.67~millions pour les seuls taux et 14.01~millions pour le seul change. Tous les facteurs sont des taux et des changes des grandes devises
avec un \omterm{def:rc:regulatory-capital-for-trading-books:ima}{horizon de liquidité} de dix jours, de sorte qu'aucune mise à l'échelle ne s'applique. L'IMCC vaut
$0.5\times14.41+0.5\times(5.67+14.01) = \text{USD}~17.04$ millions de dollars et, avec le multiplicateur de 1.5, l'exigence
est de 25.56~millions de dollars, soit 37~\% de l'exigence standardisée.
\end{example}

\begin{omfigure}
\begin{tikzpicture}
\begin{axis}[width=11.5cm,height=5.2cm,
  xlabel={fin de la fenêtre de douze mois},ylabel={ES à 97.5~\% sur dix jours (millions de dollars)},
  tick label style={font=\small},label style={font=\small},
  xmin=2017,xmax=2026.8,ymin=0,ymax=16,grid=major,grid style={gray!25},
  x tick label style={/pgf/number format/1000 sep={}},
  table/col sep=comma]
\addplot[omDef,very thick] table[x=end,y=es]{figdata/rates-credit-risk/23-regulatory-capital-for-trading-books/es_windows.csv};
\end{axis}
\end{tikzpicture}

\omcaption{L'\omterm{def:rc:market-risk-measures:es}{expected shortfall} du portefeuille actuel sur chaque fenêtre de douze mois de l'historique depuis
2016, selon la date de fin de la fenêtre. Le calibrage en période de tensions retient le maximum, soit l'année
jusqu'en octobre 2022. Données~: taux au pair du Trésor américain, taux de référence de la BCE~; tutoriel du chapitre.}\label{fig:rc:regulatory-capital-for-trading-books:windows}
\end{omfigure}

\section{Backtesting et test d'attribution du P\&L}

\begin{definition}[P\&L hypothétique et P\&L théorique du risque]\label{def:rc:regulatory-capital-for-trading-books:hpl}
Le \emph{P\&L hypothétique}\index{P\&L hypothétique} (HPL) d'un desk est la variation quotidienne de la
valeur des positions de la veille calculée par les systèmes de valorisation de la banque avec les données de marché du jour.
Le \emph{P\&L théorique du risque}\index{P\&L théorique du risque} (RTPL) est la même variation calculée par
le modèle de risque, à partir de ses propres \omterm{def:rc:market-risk-measures:factor}{facteurs de risque} et de sa propre réévaluation.
\end{definition}

\begin{definition}[Test d'attribution du P\&L]\label{def:rc:regulatory-capital-for-trading-books:pla}
Le \emph{test d'attribution du P\&L}\index{test d'attribution du P\&L} compare le HPL et le RTPL d'un desk sur
250 jours, au moyen de la corrélation de rang de Spearman de leurs rangs et de la distance de Kolmogorov--Smirnov entre
leurs distributions~: vert si la corrélation dépasse 0.80 et la distance est inférieure à 0.09, rouge si
la corrélation est inférieure à 0.70 ou la distance supérieure à 0.12, orange sinon. Les desks rouges utilisent
l'approche standardisée~; les desks orange acquittent une majoration égale à la moitié de l'écart entre l'exigence standardisée et
celle des modèles internes (pour une banque dont le seul desk est orange).
\end{definition}

\omcode{code/firm/frtb/firm_frtb.py}{136}{159}{Corrélation de rang de Spearman, distance de Kolmogorov--Smirnov et zones d'attribution.}\label{lst:rc:regulatory-capital-for-trading-books:pla}

\begin{example}[Trois modèles de risque pour un même desk]\label{ex:rc:regulatory-capital-for-trading-books:pla}
Sur les 250 derniers jours, un modèle de risque qui réévalue le portefeuille par le développement delta--gamma
passe sans difficulté (Spearman 1.00, KS 0.008). Un modèle sans facteur yen passe aussi (0.956 et 0.060).
Un modèle sans facteur à deux ans est orange par sa distribution (0.949 et 0.100). Un modèle qui
remplace le rendement à dix ans par celui à deux ans, réponse courante à un facteur jugé non modélisable,
est orange par sa corrélation~: 0.703 et 0.044 (la \cref{fig:rc:regulatory-capital-for-trading-books:pla}).
\end{example}

\begin{omfigure}
\begin{tikzpicture}
\begin{axis}[width=8cm,height=6.5cm,
  xlabel={P\&L hypothétique (millions de dollars)},ylabel={P\&L théorique du risque (millions de dollars)},
  tick label style={font=\small},label style={font=\small},
  xmin=-2.5,xmax=2,ymin=-2.5,ymax=2,grid=major,grid style={gray!25},
  table/col sep=comma]
\addplot[omExo,dashed,domain=-2.5:2,samples=2,forget plot] {x};
\addplot[omThm,only marks,mark=*,mark size=0.9pt] table[x=hpl,y=rtpl]{figdata/rates-credit-risk/23-regulatory-capital-for-trading-books/pla.csv};
\end{axis}
\end{tikzpicture}

\omcaption{\omterm{def:rc:regulatory-capital-for-trading-books:hpl}{P\&L hypothétique} quotidien du portefeuille face au P\&L d'un modèle de risque qui remplace le
rendement à dix ans par celui à deux ans, sur 250 jours. Les rangs ne concordent que faiblement (Spearman 0.703)~:
orange. Données~: Trésor américain, BCE~; tutoriel du chapitre.}\label{fig:rc:regulatory-capital-for-trading-books:pla}
\end{omfigure}

\section{Facteurs de risque non modélisables}

\begin{definition}[Test d'éligibilité des facteurs de risque, facteur de risque non modélisable]\label{def:rc:regulatory-capital-for-trading-books:nmrf}
Le \emph{test d'éligibilité des facteurs de risque}\index{test d'éligibilité des facteurs de risque} admet un \omterm{def:rc:market-risk-measures:factor}{facteur de risque}
dans le modèle interne s'il compte au moins 24 observations de prix réels dans l'année écoulée, sans
aucune période de 90 jours n'en comptant moins de quatre (ou au moins 100 dans l'année). Un facteur qui échoue est un
\emph{facteur de risque non modélisable}\index{facteur de risque non modélisable}~: il est capitalisé séparément
par son propre scénario de tensions, avec peu de diversification.
\end{definition}

\begin{dated}{2026-09}{Quand la FRTB s'applique}\label{dat:rc:regulatory-capital-for-trading-books:dates}
La norme du Comité de Bâle de janvier 2019 devait s'appliquer à partir du 1~janvier 2022~; en mars 2020, son
organe de surveillance l'a reportée au 1~janvier 2023. Dans l'Union européenne, après que les règles ont été
reportées de deux ans, elles s'appliquent à partir du 1~janvier 2027, avec des ajustements ciblés et temporaires pendant trois
ans, dont un multiplicateur de capital (acte délégué de la Commission adopté le 4 juin 2026, publié le
1~septembre 2026). Au Royaume-Uni, la PRA a consulté en juin 2026 sur ses règles de modèles internes,
applicables à partir du 1~janvier 2028, avec une période de suivi du test d'attribution étendue de un à trois
ans~; ses autres règles de Bâle~3.1 s'appliquent à partir de janvier 2027. Aux États-Unis, les agences ont proposé
en mars 2026 un cadre transposant les éléments finaux de Bâle~III, avec des règles de risque de marché pour les banques
dont l'activité de négociation est importante~; les commentaires ont été clos le 18 juin 2026.
\end{dated}

\section{Tutoriel~: deux montants de capital pour un même desk}\label{tut:rc:regulatory-capital-for-trading-books:tutorial}

\begin{tutorial}
\textbf{Objectif.} Calculer les exigences standardisée et des modèles internes du portefeuille du chapitre~21 et lancer
le test d'attribution sur plusieurs modèles de risque. \textbf{État final~:} les chiffres des
\cref{ex:rc:regulatory-capital-for-trading-books:sa,ex:rc:regulatory-capital-for-trading-books:ima,ex:rc:regulatory-capital-for-trading-books:pla}
et les trois graphiques.
\begin{tutsteps}
\item \textbf{Sensibilités}~: \texttt{sensitivities()} soumet le portefeuille aux chocs prescrits par la norme.
\item \textbf{Approche standardisée}~: \texttt{standardised()} sous les trois scénarios.
\item \textbf{Modèles internes}~: \texttt{internal\_models()} avec la fenêtre de tensions.
\item \textbf{Attribution}~: \texttt{pla(variant)}~; \texttt{fig\_rc\_frtb.py} écrit les graphiques.
\end{tutsteps}
\textbf{Pour aller plus loin.} Diviser le straddle par deux et voir quelle exigence varie le plus~; donner au
yen un \omterm{def:rc:regulatory-capital-for-trading-books:ima}{horizon de liquidité} de vingt jours et appliquer la formule de mise à l'échelle.
\end{tutorial}

\section{Construction~: le moteur FRTB}\label{bld:rc:regulatory-capital-for-trading-books:frtb}

\begin{build}
\bfield{Rôle} Le capital réglementaire au titre du risque de marché de la banque selon les deux approches et ses
tests d'éligibilité des desks, alimentés par le moteur de risque du chapitre~29.
\bfield{Interface} \texttt{girr\_rho}, \texttt{scenario}, \texttt{within}, \texttt{across}~;
\texttt{girr\_delta}, \texttt{fx\_delta}, \texttt{fx\_vega}, \texttt{cvr}, \texttt{fx\_curvature},
\texttt{sbm}~; \texttt{es}, \texttt{es\_liquidity}, \texttt{imcc}, \texttt{ima\_capital},
\texttt{amber\_surcharge}~; \texttt{spearman}, \texttt{ks\_stat}, \texttt{pla\_zone}~;
\texttt{rfet}.
\bfield{Règles} Paramètres conformes à la norme de Bâle de janvier 2019~; classes de risque GIRR et de change
uniquement~; devises spécifiées avec la réduction en $\sqrt2$~; le total du scénario le plus défavorable.
\bfield{Tests d'acceptation} \texttt{code/firm/frtb/tests/}~: l'exemple de corrélation de la norme
(88.69~\%)~; transformations de scénarios~; agrégation avec son repli~; signes de la courbure~; formule du capital
et majoration~; mise à l'échelle de la liquidité~; statistiques d'attribution et bornes des zones~; test d'éligibilité.
\bfield{Pour aller plus loin} Spread de crédit, classes d'actions et de matières premières~; l'\omterm{def:rc:regulatory-capital-for-trading-books:drc}{exigence au titre du risque de défaut}~; la
\omterm{def:rc:regulatory-capital-for-trading-books:drc}{majoration pour risque résiduel}~; scénarios de tensions non modélisables~; variantes nationales (multiplicateurs de l'UE).
\end{build}

\begin{omsources}
\item Basel Committee on Banking Supervision, \emph{Minimum capital requirements for market
risk}, janvier 2019.
\item Bank for International Settlements, communiqué de presse du 27 mars 2020 sur le report de Bâle~III.
\item European Commission, «~Commission adopts temporary adjustments to Basel~III market risk
rules~», 4 juin 2026.
\item Bank of England, «~PRA sets out adjustments to its market risk internal model approach
under Basel~3.1~», 19 juin 2026.
\item Board of Governors of the Federal Reserve System, communiqué de presse du 19 mars 2026 sur
des propositions visant à moderniser le cadre de fonds propres réglementaires.
\end{omsources}

\section{Exercices}

\begin{exercise}[$\star$]\label{exo:rc:regulatory-capital-for-trading-books:1}
Calculer la corrélation GIRR entre les maturités de deux et dix ans et ses valeurs dans les scénarios
de corrélation haute et basse.
\end{exercise}

\begin{exercise}[$\star$]\label{exo:rc:regulatory-capital-for-trading-books:2}
Pourquoi le scénario de corrélation basse a-t-il donné le total le plus élevé pour ce portefeuille~?
\end{exercise}

\begin{exercise}[$\star$]\label{exo:rc:regulatory-capital-for-trading-books:3}
Les statistiques d'attribution d'un desk valent 0.82 et 0.10. Dans quelle zone se trouve-t-il, et que se passe-t-il~?
\end{exercise}

\begin{exercise}[$\star\star$]\label{exo:rc:regulatory-capital-for-trading-books:4}
Pourquoi l'exigence des modèles internes est-elle nettement inférieure à celle de l'approche standardisée pour ce portefeuille~?
\end{exercise}

\begin{exercise}[$\star\star$]\label{exo:rc:regulatory-capital-for-trading-books:5}
Un \omterm{def:rc:market-risk-measures:factor}{facteur de risque} compte 30 observations dans l'année, toutes situées dans ses quatre premiers mois. Est-il modélisable~?
\end{exercise}

\begin{exercise}[$\star\star$]\label{exo:rc:regulatory-capital-for-trading-books:6}
Pourquoi l'IMCC combine-t-il l'ES diversifié avec la somme des ES des classes de risque~?
\end{exercise}

\begin{exercise}[$\star\star\star$]\label{exo:rc:regulatory-capital-for-trading-books:7}
\emph{Programmation.} Calculer l'ES d'un portefeuille dont les \omterm{def:rc:market-risk-measures:factor}{facteurs de risque} se répartissent en un groupe à dix jours d'ES~5
et un groupe à vingt jours d'ES~3 (le second étant aussi compté dans le premier), en utilisant la formule de l'\omterm{def:rc:regulatory-capital-for-trading-books:ima}{horizon
de liquidité}.
\end{exercise}

\begin{exercise}[$\star\star\star$]\label{exo:rc:regulatory-capital-for-trading-books:8}
\emph{Trouver la faille.} «~Notre modèle interne donne un tiers du capital standardisé~; le modèle est
meilleur, donc le régulateur devrait nous laisser l'utiliser pour chaque desk.~»
\end{exercise}

\section{Problème~: le desk qui a échoué au test d'attribution}

\begin{problem}[{Problème du week-end --- ambre}]\label{pb:rc:regulatory-capital-for-trading-books:1}
Le portefeuille du chapitre~21 est le seul desk d'une banque recourant à l'\omterm{def:rc:regulatory-capital-for-trading-books:ima}{approche des modèles internes}. Son modèle de risque
ne dispose pas de données fiables sur dix ans et approxime le rendement à dix ans par celui à deux ans.

\medskip\noindent\textbf{Partie I --- Deux chiffres.}
\begin{enumerate}
\item Donner l'exigence standardisée et le scénario qui la fixe.
\item Donner l'exigence des modèles internes et ses composantes.
\item Quelle composante pilote l'exigence standardisée~?
\item Pourquoi la fenêtre de tensions est-elle l'année jusqu'à octobre 2022 pour ce portefeuille~?
\item Que changerait un \omterm{def:rc:regulatory-capital-for-trading-books:ima}{horizon de liquidité} plus long pour la volatilité du change~?
\end{enumerate}

\medskip\noindent\textbf{Partie II --- Le test.}
\begin{enumerate}[resume]
\item Donner la corrélation de rang de Spearman et la distance de KS du modèle de substitution.
\item Donner la zone et la raison.
\item Donner la majoration et le capital qui en résulte.
\item Que signifierait le rouge pour le capital du desk~?
\item Lesquels des autres modèles de risque réussiraient, et lesquels échoueraient~?
\end{enumerate}

\medskip\noindent\textbf{Partie III --- Le corriger.}
\begin{enumerate}[resume]
\item Pourquoi le modèle de substitution échoue-t-il au test de corrélation de rang mais pas au test de distribution~?
\item Quelles données rendraient modélisable le rendement à dix ans~?
\item Un meilleur modèle de substitution (une régression sur le rendement à deux ans) aiderait-il~?
\item Que gagne un desk à une période de suivi plus longue, comme le propose la PRA~?
\item Comment la direction métier du desk devrait-elle trancher entre la correction et la majoration~?
\end{enumerate}

\medskip\noindent\textbf{Partie IV --- Jugement.}
\begin{enumerate}[resume]
\item Pourquoi tester tout de même l'attribution alors que le backtesting existe~?
\item Quelles incitations la majoration ambre crée-t-elle~?
\item Comment les multiplicateurs de l'UE et le report britannique modifient-ils les choix d'une banque mondiale~?
\item Énoncer le \emph{résultat nommé}~: le capital du desk en modèles internes et en approche standardisée, ainsi que les
valeurs de Spearman et de KS qui le placent en ambre.
\item En une phrase~: que vérifie réellement le test d'attribution~?
\end{enumerate}
\end{problem}

\section{Questions d'entretien}

\begin{interviewq}[$\star$ \iqroles{risque, banque}]\label{iq:rc:regulatory-capital-for-trading-books:1}
Qu'est-ce qui a changé de Bâle~2.5 à la FRTB~?
\end{interviewq}

\begin{interviewq}[$\star\star$ \iqroles{risque, développeur}]\label{iq:rc:regulatory-capital-for-trading-books:2}
Comment appliqueriez-vous la \omterm{def:rc:regulatory-capital-for-trading-books:sbm}{méthode fondée sur les sensibilités} à un portefeuille de swaps de taux~?
\end{interviewq}

\begin{interviewq}[$\star\star$ \iqroles{risque}]\label{iq:rc:regulatory-capital-for-trading-books:3}
Qu'est-ce que le \omterm{def:rc:regulatory-capital-for-trading-books:pla}{test d'attribution du P\&L} et pourquoi les desks y échouent-ils~?
\end{interviewq}

\begin{interviewq}[$\star\star$ \iqroles{chercheur, risque}]\label{iq:rc:regulatory-capital-for-trading-books:4}
Pourquoi l'\omterm{def:rc:market-risk-measures:es}{expected shortfall} avec \omterm{def:rc:regulatory-capital-for-trading-books:ima}{horizons de liquidité} plutôt qu'une VaR à dix jours~?
\end{interviewq}

\begin{interviewq}[$\star\star\star$ \iqroles{banque}]\label{iq:rc:regulatory-capital-for-trading-books:5}
Comment décideriez-vous quels desks doivent demander l'agrément des modèles internes~?
\end{interviewq}

\begin{interviewq}[$\star\star\star$ \iqroles{développeur}]\label{iq:rc:regulatory-capital-for-trading-books:6}
De quelles données et de quels systèmes une banque a-t-elle besoin pour appliquer l'\omterm{def:rc:regulatory-capital-for-trading-books:ima}{approche des modèles internes}~?
\end{interviewq}
